% Conversão fiel de chapter-04-text-h-0044; origem RELATORIO_ATUAL.odt.
\hypertarget{__RefHeading___Toc4777_3118176633}{}
\refstepcounter{subsection}
\subsection*{4.2.3 Arquitetura utilizada}
\label{ages-ii:text-h-0044}
\addcontentsline{toc}{subsection}{4.2.3 Arquitetura utilizada}

\href{https://tools.ages.pucrs.br/industria-digital/wiki/-/wikis/Arquitetura}{}

% SERIALIZATION text-p-0785 original_whitespace
\indent A arquitetura utilizada no projeto segue um padrão amplamente adotado, conhecido por sua separação clara de responsabilidades entre as camadas. No Front-end, o \textbf{Angular} é responsável por criar a interface de usuário interativa, permitindo que o usuário final realize ações como buscas de empresas. Essas interações enviam requisições HTTP (GET, POST, PUT, DELETE) para o Back-end, que utiliza o \textbf{FastAPI}\originalspace{} para processar as requisições. O FastAPI atua como uma ponte entre a interface e o banco de dados \textbf{PostgreSQL}, enviando consultas SQL e retornando dados relevantes. Esses dados são retornados em formato JSON para o Angular, que então atualiza a interface do usuário em tempo real. Essa arquitetura modular e escalável garante que a apresentação, a lógica de negócios e a persistência de dados sejam gerenciadas de forma independente, permitindo uma manutenção mais fácil e flexibilidade para expansões futuras, que serão necessárias.\par

% SERIALIZATION text-p-0788 abnt_figure_caption
\memorialfigure{imagens/original/asset-0008.png}{\textbf{Figura 10 –} Arquitetura da Aplicação - AGES II}{ages-ii:draw-image-0010}{Fonte: pendente de identificação}{0.92} % <<ASSET draw-image-0010 asset-0008>>

\indent \href{https://tools.ages.pucrs.br/industria-digital/wiki/-/wikis/Arquitetura}{Wiki da Arquitetura}\par
